% Options for packages loaded elsewhere
\PassOptionsToPackage{unicode}{hyperref}
\PassOptionsToPackage{hyphens}{url}
\PassOptionsToPackage{dvipsnames,svgnames*,x11names*}{xcolor}
%
\documentclass[
]{article}
\usepackage{lmodern}
\usepackage{amsmath}
\usepackage{ifxetex,ifluatex}
\ifnum 0\ifxetex 1\fi\ifluatex 1\fi=0 % if pdftex
  \usepackage[T1]{fontenc}
  \usepackage[utf8]{inputenc}
  \usepackage{textcomp} % provide euro and other symbols
  \usepackage{amssymb}
\else % if luatex or xetex
  \usepackage{unicode-math}
  \defaultfontfeatures{Scale=MatchLowercase}
  \defaultfontfeatures[\rmfamily]{Ligatures=TeX,Scale=1}
\fi
% Use upquote if available, for straight quotes in verbatim environments
\IfFileExists{upquote.sty}{\usepackage{upquote}}{}
\IfFileExists{microtype.sty}{% use microtype if available
  \usepackage[]{microtype}
  \UseMicrotypeSet[protrusion]{basicmath} % disable protrusion for tt fonts
}{}
\makeatletter
\@ifundefined{KOMAClassName}{% if non-KOMA class
  \IfFileExists{parskip.sty}{%
    \usepackage{parskip}
  }{% else
    \setlength{\parindent}{0pt}
    \setlength{\parskip}{6pt plus 2pt minus 1pt}}
}{% if KOMA class
  \KOMAoptions{parskip=half}}
\makeatother
\usepackage{xcolor}
\IfFileExists{xurl.sty}{\usepackage{xurl}}{} % add URL line breaks if available
\IfFileExists{bookmark.sty}{\usepackage{bookmark}}{\usepackage{hyperref}}
\hypersetup{
  pdftitle={Taller 1 entrega problema en grupo. MAT3 (estadística) GIN2 2020-2021 - Probabilidad, Variables Aleatorias, Distribuciones Notables 28-03-2020.},
  pdfauthor={Marc Cañellas, Eduardo Bonnin, Jaume Adrover, Joan Balaguer},
  colorlinks=true,
  linkcolor=Maroon,
  filecolor=Maroon,
  citecolor=Blue,
  urlcolor=blue,
  pdfcreator={LaTeX via pandoc}}
\urlstyle{same} % disable monospaced font for URLs
\usepackage[margin=1in]{geometry}
\usepackage{color}
\usepackage{fancyvrb}
\newcommand{\VerbBar}{|}
\newcommand{\VERB}{\Verb[commandchars=\\\{\}]}
\DefineVerbatimEnvironment{Highlighting}{Verbatim}{commandchars=\\\{\}}
% Add ',fontsize=\small' for more characters per line
\usepackage{framed}
\definecolor{shadecolor}{RGB}{248,248,248}
\newenvironment{Shaded}{\begin{snugshade}}{\end{snugshade}}
\newcommand{\AlertTok}[1]{\textcolor[rgb]{0.94,0.16,0.16}{#1}}
\newcommand{\AnnotationTok}[1]{\textcolor[rgb]{0.56,0.35,0.01}{\textbf{\textit{#1}}}}
\newcommand{\AttributeTok}[1]{\textcolor[rgb]{0.77,0.63,0.00}{#1}}
\newcommand{\BaseNTok}[1]{\textcolor[rgb]{0.00,0.00,0.81}{#1}}
\newcommand{\BuiltInTok}[1]{#1}
\newcommand{\CharTok}[1]{\textcolor[rgb]{0.31,0.60,0.02}{#1}}
\newcommand{\CommentTok}[1]{\textcolor[rgb]{0.56,0.35,0.01}{\textit{#1}}}
\newcommand{\CommentVarTok}[1]{\textcolor[rgb]{0.56,0.35,0.01}{\textbf{\textit{#1}}}}
\newcommand{\ConstantTok}[1]{\textcolor[rgb]{0.00,0.00,0.00}{#1}}
\newcommand{\ControlFlowTok}[1]{\textcolor[rgb]{0.13,0.29,0.53}{\textbf{#1}}}
\newcommand{\DataTypeTok}[1]{\textcolor[rgb]{0.13,0.29,0.53}{#1}}
\newcommand{\DecValTok}[1]{\textcolor[rgb]{0.00,0.00,0.81}{#1}}
\newcommand{\DocumentationTok}[1]{\textcolor[rgb]{0.56,0.35,0.01}{\textbf{\textit{#1}}}}
\newcommand{\ErrorTok}[1]{\textcolor[rgb]{0.64,0.00,0.00}{\textbf{#1}}}
\newcommand{\ExtensionTok}[1]{#1}
\newcommand{\FloatTok}[1]{\textcolor[rgb]{0.00,0.00,0.81}{#1}}
\newcommand{\FunctionTok}[1]{\textcolor[rgb]{0.00,0.00,0.00}{#1}}
\newcommand{\ImportTok}[1]{#1}
\newcommand{\InformationTok}[1]{\textcolor[rgb]{0.56,0.35,0.01}{\textbf{\textit{#1}}}}
\newcommand{\KeywordTok}[1]{\textcolor[rgb]{0.13,0.29,0.53}{\textbf{#1}}}
\newcommand{\NormalTok}[1]{#1}
\newcommand{\OperatorTok}[1]{\textcolor[rgb]{0.81,0.36,0.00}{\textbf{#1}}}
\newcommand{\OtherTok}[1]{\textcolor[rgb]{0.56,0.35,0.01}{#1}}
\newcommand{\PreprocessorTok}[1]{\textcolor[rgb]{0.56,0.35,0.01}{\textit{#1}}}
\newcommand{\RegionMarkerTok}[1]{#1}
\newcommand{\SpecialCharTok}[1]{\textcolor[rgb]{0.00,0.00,0.00}{#1}}
\newcommand{\SpecialStringTok}[1]{\textcolor[rgb]{0.31,0.60,0.02}{#1}}
\newcommand{\StringTok}[1]{\textcolor[rgb]{0.31,0.60,0.02}{#1}}
\newcommand{\VariableTok}[1]{\textcolor[rgb]{0.00,0.00,0.00}{#1}}
\newcommand{\VerbatimStringTok}[1]{\textcolor[rgb]{0.31,0.60,0.02}{#1}}
\newcommand{\WarningTok}[1]{\textcolor[rgb]{0.56,0.35,0.01}{\textbf{\textit{#1}}}}
\usepackage{graphicx}
\makeatletter
\def\maxwidth{\ifdim\Gin@nat@width>\linewidth\linewidth\else\Gin@nat@width\fi}
\def\maxheight{\ifdim\Gin@nat@height>\textheight\textheight\else\Gin@nat@height\fi}
\makeatother
% Scale images if necessary, so that they will not overflow the page
% margins by default, and it is still possible to overwrite the defaults
% using explicit options in \includegraphics[width, height, ...]{}
\setkeys{Gin}{width=\maxwidth,height=\maxheight,keepaspectratio}
% Set default figure placement to htbp
\makeatletter
\def\fps@figure{htbp}
\makeatother
\setlength{\emergencystretch}{3em} % prevent overfull lines
\providecommand{\tightlist}{%
  \setlength{\itemsep}{0pt}\setlength{\parskip}{0pt}}
\setcounter{secnumdepth}{-\maxdimen} % remove section numbering
\renewcommand{\contentsname}{Contenidos}
\ifluatex
  \usepackage{selnolig}  % disable illegal ligatures
\fi

\title{Taller 1 entrega problema en grupo. MAT3 (estadística) GIN2
2020-2021 - Probabilidad, Variables Aleatorias, Distribuciones Notables
28-03-2020.}
\author{Marc Cañellas, Eduardo Bonnin, Jaume Adrover, Joan Balaguer}
\date{}

\begin{document}
\maketitle

\hypertarget{taller1-evaluable.-entrega-de-problemas}{%
\section{Taller1 evaluable. Entrega de
problemas}\label{taller1-evaluable.-entrega-de-problemas}}

Taller en grupo entregad las soluciones en .Rmd y .html o .pdf. o
escribidlas de forma manual y escanear el resultado, en un solo fichero.

\hypertarget{problema-1}{%
\subsection{Problema 1}\label{problema-1}}

Sean \(A\) , \(B\) y \(C\) tres sucesos tales que \(P(A)=0.4\),
\(P(B)=0.4\) y \(P(A\cup B)=0.6\). Calcular \(P(A\cap B)\).

\hypertarget{soluciuxf3n-ejercicio-1}{%
\subsection{Solución Ejercicio 1}\label{soluciuxf3n-ejercicio-1}}

Sabiendo que la \(P(A\cap B)= P(A)+P(B)-P(A\cup B)\) y sabiendo el valor
de \(P(A)=0.4\), \(P(B)=0.4\) y \(P(A\cup B)=0.6\), entonces
\(\longrightarrow\) \(P(A\cap B)= 0.4+0.4-0.6 = 0.2\).

\hypertarget{soluciuxf3n-con-r-ejercicio-1}{%
\subsection{\texorpdfstring{Solución con \texttt{R} Ejercicio
1}{Solución con R Ejercicio 1}}\label{soluciuxf3n-con-r-ejercicio-1}}

\begin{Shaded}
\begin{Highlighting}[]
\NormalTok{A}\OtherTok{=}\FloatTok{0.4}
\NormalTok{B}\OtherTok{=}\FloatTok{0.4}
\NormalTok{AunionB}\OtherTok{=}\FloatTok{0.6}
\NormalTok{AintersecB}\OtherTok{=}\NormalTok{A}\SpecialCharTok{+}\NormalTok{B}\SpecialCharTok{{-}}\NormalTok{AunionB}
\NormalTok{AintersecB}
\end{Highlighting}
\end{Shaded}

\begin{verbatim}
## [1] 0.2
\end{verbatim}

\hypertarget{problema-2}{%
\subsection{Problema 2}\label{problema-2}}

Consideremos la v.a. continua \(X\) que tiene por función de densidad
para a alguna constante \(\alpha\in \mathbb{R}\):

\[
f_X (t)=
\left\{\begin{array}{ll}
\alpha \cdot t^4, & \mbox{si } -1 < t <1,
 \\
0 & \mbox{ en otro caso}.
\end{array}\right.
\]

\begin{enumerate}
\def\labelenumi{\arabic{enumi}.}
\tightlist
\item
  Calculad \(\alpha\) para que \(f_X\) sea densidad y especificad su
  dominio \(D_X\).
\item
  Calculad la función de distribución de la v.a. \(X\);
  \(F_X(x)=P(X\leq x)\).
\item
  Calculad \(E(X)\) y \(Var(X)\).
\item
  Calcula en cuantil \(0.9\) de \(X\).
\end{enumerate}

\hypertarget{soluciuxf3n-ejercicio-2}{%
\subsection{Solución Ejercicio 2}\label{soluciuxf3n-ejercicio-2}}

\begin{enumerate}
\def\labelenumi{\arabic{enumi}.}
\item
  Sabemos que una de las propiedades de toda funcion de densidad es que
  \(\int_{-\infty}^{+\infty}f_{X}(t)=1\). Por lo tanto, sabiendo que
  \(D_{X}=t\in (-1,1)\): \(\\\) \(\\\)
  \[\int_{-1}^{+1} \alpha \cdot t^4 = 1 \mbox{ }\Rightarrow \mbox{ } [\frac{\alpha \cdot t^5}{5}]_{-1}^{+1} = 1 \mbox{ }\Rightarrow \mbox{ } \alpha \cdot (\frac{1}{5}+\frac{1}{5})=1  \mbox{ }\Rightarrow \mbox{ } \alpha= \frac{5}{2}\]
  \(\\\) \(\\\)
\item
  Para calcular la función de distribución, tenemos que tener en cuenta
  que: dada una V.A continua \(X\),
  \(F_{X}(t)=\int_{-\infty}^{x}f_{X}(t)\), siendo \(F_{X}(t)\) su
  función de distribución y \(f_{X}(t)\) su funcion de densidad. Por lo
  tanto: \(\\\) \(\\\)
  \[\int_{-1}^{x} f_{X}(t)=\int_{-1}^{x} \alpha \cdot t^4 \mbox{ } dt= [\frac {t^5}{2}]_{-1}^{x}= \frac {x^5}{2}-(-\frac{1}{2})= \frac{x^5+1}{2}\]
  \(\\\) \(\\\) Finalmente, la función de distribución de la V.A \(X\)
  queda de la siguiente forma: \(\\\) \(\\\) \[F_{X}(t)=
  \left\{\begin{array}{ll}
  0 &\mbox{si }x\le-1 
  \\
  \frac{x^5+1}{2}&\mbox{si }-1<x<1
  \\
  1&\mbox{si }x\ge1
  \end{array}\right.\] \(\\\)
\item
  Sabemos que
  \(E(t)=\int_{-\infty}^{+\infty} t\cdot f_{X}(t)\mbox{ }dt\), y si
  integramos en el dominio del ejercicio: \(\\\) \(\\\)
  \[ E(t)=\int_{-1}^{+1} t \cdot \alpha \cdot t^4 \mbox{ }dt= [\frac{\alpha \cdot t^6}{6}]_{-1}^{+1}= \frac{\alpha}{6}- \frac{\alpha}{6}=0\]
  \(\\\) \(\\\) Para calcular la varianza, usaremos la siguiente
  expressión \(\longrightarrow Var(X)=E((t-\mu_{t})^{2})\), siendo
  \(\mu_{t}\) la esperanza de la V.A \(X\). En cosecuencia el valor de
  \(Var(X)\) es el siguiente: \(\\\) \(\\\)
  \[Var(X)=E(t^2)=\int_{-1}^{1}t^6 \cdot \alpha \mbox{ }dt= [\frac{\alpha \cdot t^7}{7}]_{-1}^{+1}= \frac{\alpha}{7}+\frac{\alpha}{7}= \frac{2 \cdot \alpha}{7}= \frac{5}{7}= 0.7142857\]
  \(\\\)
\item
  Por definicion, un cuantil es \(P(X \le t_{p}) \ge p\) es decir, es el
  primer valor de \(t \in D_{X}\) tal que supera por primera vez a
  p.~Como estamos trabajando con una V.A continua, deberemos usar la
  propiedad que nos indica que \(F_{X}^{-1}(p)=t_{p}\). Por tanto, vamos
  a calcular la inversa de la funcion de distribución: \(\\\) \(\\\)
  \[F_{X}^{-1}(t)=
  \left\{\begin{array}{ll}
  \sqrt[5]{2\cdot t}, & \mbox{si } -1 < t <1,
   \\
  0 & \mbox{ en otro caso}.
  \end{array}\right.\] \(\\\) \(\\\) En el ejercicio nos piden que
  calculemos el cuantil \(0,9\) de \(X\): \(\\\) \(\\\)
  \[t_{0.9}=F_{X}^{-1}(0.9)= \sqrt[5]{2 \cdot 0.9}= 1.124746\]
\end{enumerate}

\hypertarget{soluciuxf3n-con-r-ejercicio-2}{%
\subsection{\texorpdfstring{Solución con \texttt{R} Ejercicio
2}{Solución con R Ejercicio 2}}\label{soluciuxf3n-con-r-ejercicio-2}}

\begin{Shaded}
\begin{Highlighting}[]
\CommentTok{\# Apartado 1}

\NormalTok{alpha}\OtherTok{=}\DecValTok{5}\SpecialCharTok{/}\DecValTok{2}
\NormalTok{alpha}
\end{Highlighting}
\end{Shaded}

\begin{verbatim}
## [1] 2.5
\end{verbatim}

\begin{Shaded}
\begin{Highlighting}[]
\CommentTok{\# Apartado 3}

\NormalTok{Esperanza1}\OtherTok{=}\NormalTok{(alpha}\SpecialCharTok{/}\DecValTok{6}\NormalTok{)}\SpecialCharTok{{-}}\NormalTok{(alpha}\SpecialCharTok{/}\DecValTok{6}\NormalTok{)}
\NormalTok{Esperanza1}
\end{Highlighting}
\end{Shaded}

\begin{verbatim}
## [1] 0
\end{verbatim}

\begin{Shaded}
\begin{Highlighting}[]
\NormalTok{Varianza1}\OtherTok{=}\NormalTok{(alpha}\SpecialCharTok{/}\DecValTok{7}\NormalTok{)}\SpecialCharTok{+}\NormalTok{(alpha}\SpecialCharTok{/}\DecValTok{7}\NormalTok{)}
\NormalTok{Varianza1}
\end{Highlighting}
\end{Shaded}

\begin{verbatim}
## [1] 0.7142857
\end{verbatim}

\begin{Shaded}
\begin{Highlighting}[]
\CommentTok{\# Apartado 4}

\NormalTok{Cuantil1}\OtherTok{=}\NormalTok{ (}\DecValTok{2}\SpecialCharTok{*}\FloatTok{0.9}\NormalTok{)}\SpecialCharTok{\^{}}\NormalTok{(}\DecValTok{1}\SpecialCharTok{/}\DecValTok{5}\NormalTok{)}
\NormalTok{Cuantil1}
\end{Highlighting}
\end{Shaded}

\begin{verbatim}
## [1] 1.124746
\end{verbatim}

\hypertarget{problema-3}{%
\subsection{Problema 3}\label{problema-3}}

Sea \(Y\) una variable discreta con función de probabilidad :

\[
P_Y(y)=
\left\{\begin{array}{ll}
\alpha\cdot\frac{1}{x^2} & \mbox{,si } x=-2,-1,1,2,
 \\
0 & \mbox{ en otro caso}.
\end{array}\right. 
\] 1. Hallad \(\alpha\) para que \(P_Y\) sea función de probabilidad. 2.
Hallad la función de distribución \(F_Y(y)=P(Y\leq Y)\). 3. Calculad
\(E(Y)\) y \(Var(Y)\) 4. Calculad el cuantil 0.5 de de \(Y\)

\hypertarget{soluciuxf3n-ejercicio-3}{%
\subsection{Solución Ejercicio 3}\label{soluciuxf3n-ejercicio-3}}

\begin{enumerate}
\def\labelenumi{\arabic{enumi}.}
\tightlist
\item
  Sabemos que para toda V.A discreta, su función de probabilidad
  cumplirá la siguente propiedad: \(\sum P(y)=1\) para toda
  \(y \in D_{Y}\). Sabemos también que
  \(D_{Y} = y \in \mbox{ } [-2, -1,\mbox{ } 1,\mbox{ } 2]\) y por lo
  tanto: \(\\\) \(\\\)
  \[P(-2)+P(-1)+P(1)+P(2)=1 \mbox{ }\Rightarrow \mbox{ } \frac{\alpha}{4}+\frac{\alpha}{1}+\frac{\alpha}{1}+\frac{\alpha}{4}=1 \mbox{ }\Rightarrow \mbox{ } 2\cdot \alpha \cdot (\frac{5}{4})=1 \mbox{ }\Rightarrow \mbox{ } \alpha = \frac{2}{5}= 0.4\]
  \(\\\) \(\\\)
\item
  La función de distribución de la V.A Y será la siguente: \(\\\) \(\\\)
  \[F_{Y}(y)= \left\{ \begin{array}{lcc}
            P(-2) &   si  & y=-2
            \\ P(-2)+P(-1) &  si & y=-1
            \\ P(-2)+P(-1)+P(1) &  si  & y = \mbox{ }1
            \\ P(-2)+P(-1)+P(1)+P(2) & si & y= \mbox{ }2
            \end{array}
  \right.\]
\end{enumerate}

Sabiendo que: \(\\\) \(\\\)
\[P(-2)= \frac{\alpha}{4} = \frac{0.4}{4} = \frac{1}{10}=0.1\]
\[P(-1)= \frac{\alpha}{1} = \frac{2}{5} = 0.4\]
\[P(1)= \frac{\alpha}{1} = \frac{2}{5} = 0.4\]
\[P(2)= \frac{\alpha}{4} = \frac{0.4}{4} = \frac{1}{10}=0.1\] \(\\\)
\(\\\) Finalmente nos quedará la siguente función de distribución:
\(\\\) \(\\\) \[F_{Y}(y)= \left\{ \begin{array}{lcc}
              0 & si & y < 2
             \\0.1 &   si  & y=-2
             \\ 0.5 &  si & y=-1
             \\ 0.9 &  si  & y = \mbox{ }1
             \\ 1 & si & y \ge \mbox{ } 2
             \\
             \end{array}
   \right.\] \(\\\) \(\\\) 3.
\(E(Y)= \sum y \cdot P_{Y}(y)= -2\cdot P(-2)-1\cdot P(-1)+1\cdot P(1)+2\cdot P(2)=-2 \cdot 0.1-1 \cdot 0.4+1 \cdot 0.4+2 \cdot 0.1=0\)
\(\\\) \(\\\)
\(Var(Y)= E((y-E(Y))^2)= E(y^2) = \sum y^2 \cdot P_{Y}(y) =4 \cdot 0.1+1 \cdot 0.4+1 \cdot 0.4+4 \cdot 0.1= \frac{8}{5}=1.6\)
\(\\\) \(\\\) 4. Para calcular el cunatil de una V.A discreta, basta que
observemos qué valor de y toma la función de distribución cuando la
probabilidad resultante toma un valor superior o igual a 0.5. En este
caso, dicho valor es \(y_{0.5}=-1\) \[\\\]

\hypertarget{soluciuxf3n-con-r-ejercicio-3}{%
\subsection{\texorpdfstring{Solución con \texttt{R} Ejercicio
3}{Solución con R Ejercicio 3}}\label{soluciuxf3n-con-r-ejercicio-3}}

\begin{Shaded}
\begin{Highlighting}[]
\CommentTok{\#Apartado 1}

\NormalTok{alpha}\OtherTok{=}\NormalTok{(}\DecValTok{1}\SpecialCharTok{/}\DecValTok{2}\NormalTok{)}\SpecialCharTok{*}\NormalTok{(}\DecValTok{4}\SpecialCharTok{/}\DecValTok{5}\NormalTok{)}
\NormalTok{alpha}
\end{Highlighting}
\end{Shaded}

\begin{verbatim}
## [1] 0.4
\end{verbatim}

\begin{Shaded}
\begin{Highlighting}[]
\CommentTok{\#Apartado 2}

\NormalTok{P1}\OtherTok{=}\NormalTok{(alpha}\SpecialCharTok{/}\DecValTok{4}\NormalTok{)}
\NormalTok{P2}\OtherTok{=}\NormalTok{alpha}
\NormalTok{P3}\OtherTok{=}\NormalTok{P2}
\NormalTok{P4}\OtherTok{=}\NormalTok{P1}
\NormalTok{P1}
\end{Highlighting}
\end{Shaded}

\begin{verbatim}
## [1] 0.1
\end{verbatim}

\begin{Shaded}
\begin{Highlighting}[]
\NormalTok{P2}
\end{Highlighting}
\end{Shaded}

\begin{verbatim}
## [1] 0.4
\end{verbatim}

\begin{Shaded}
\begin{Highlighting}[]
\NormalTok{P3}
\end{Highlighting}
\end{Shaded}

\begin{verbatim}
## [1] 0.4
\end{verbatim}

\begin{Shaded}
\begin{Highlighting}[]
\NormalTok{P4}
\end{Highlighting}
\end{Shaded}

\begin{verbatim}
## [1] 0.1
\end{verbatim}

\begin{Shaded}
\begin{Highlighting}[]
\NormalTok{P1}
\end{Highlighting}
\end{Shaded}

\begin{verbatim}
## [1] 0.1
\end{verbatim}

\begin{Shaded}
\begin{Highlighting}[]
\NormalTok{P1}\SpecialCharTok{+}\NormalTok{P2}
\end{Highlighting}
\end{Shaded}

\begin{verbatim}
## [1] 0.5
\end{verbatim}

\begin{Shaded}
\begin{Highlighting}[]
\NormalTok{P1}\SpecialCharTok{+}\NormalTok{P2}\SpecialCharTok{+}\NormalTok{P3}
\end{Highlighting}
\end{Shaded}

\begin{verbatim}
## [1] 0.9
\end{verbatim}

\begin{Shaded}
\begin{Highlighting}[]
\NormalTok{P1}\SpecialCharTok{+}\NormalTok{P2}\SpecialCharTok{+}\NormalTok{P3}\SpecialCharTok{+}\NormalTok{P4}
\end{Highlighting}
\end{Shaded}

\begin{verbatim}
## [1] 1
\end{verbatim}

\begin{Shaded}
\begin{Highlighting}[]
\CommentTok{\#Apartado 3}

\NormalTok{Esperanza2}\OtherTok{=}\NormalTok{ (}\SpecialCharTok{{-}}\DecValTok{2}\NormalTok{)}\SpecialCharTok{*}\FloatTok{0.1}\DecValTok{{-}1}\SpecialCharTok{*}\FloatTok{0.4}\SpecialCharTok{+}\DecValTok{1}\SpecialCharTok{*}\FloatTok{0.4}\SpecialCharTok{+}\DecValTok{2}\SpecialCharTok{*}\FloatTok{0.1}
\CommentTok{\#Deberia de ser 0}
\NormalTok{Esperanza2}
\end{Highlighting}
\end{Shaded}

\begin{verbatim}
## [1] -5.551115e-17
\end{verbatim}

\begin{Shaded}
\begin{Highlighting}[]
\NormalTok{Varianza2}\OtherTok{=}\DecValTok{4}\SpecialCharTok{*}\FloatTok{0.1+0.4+0.4}\SpecialCharTok{+}\DecValTok{4}\SpecialCharTok{*}\FloatTok{0.1}
\NormalTok{Varianza2}
\end{Highlighting}
\end{Shaded}

\begin{verbatim}
## [1] 1.6
\end{verbatim}

\hypertarget{problema-4}{%
\subsection{Problema 4}\label{problema-4}}

Tenemos un dado, bien equilibrado, de doce caras numeradas del 1 al 12
(\href{https://es.wikipedia.org/wiki/Dados_de_rol}{dodecaedro dados de
rol}).

\begin{enumerate}
\def\labelenumi{\arabic{enumi}.}
\tightlist
\item
  Calcular la función de probabilidad de la variables \(X=\) número de
  la cara superior del dado en un lanzamiento.
\item
  Calcular la función de distribución de \(X\) y el cuantil \(0.4\).
\item
  Si \(Y\) es al v.a.que cuenta el número de veces que tiramos el dado
  hasta obtener el primer \(5\) calcular la función de distribución de
  \(Y\)
\item
  ¿Qué valor tienen \(E(X)\) y \(Var(X)\).
\end{enumerate}

\hypertarget{soluciuxf3n-ejercicio-4}{%
\subsection{Solución Ejercicio 4}\label{soluciuxf3n-ejercicio-4}}

\begin{enumerate}
\def\labelenumi{\arabic{enumi}.}
\tightlist
\item
  Como la V.A \(X\) representa el numero de la cara superior del dado en
  un lanzamiento, consideraremos el siguente espacio mostral
  \(X(\Omega)=[ 1, 2, 3, 4, 5, 6, 7, 8, 9,1 0, 11, 12 ]\).Además, como
  nos dicen que el dado está equilibrado, las probabilidades de todos
  los sucesos son equiprobables \(P(x \in X(\Omega))= \frac{1}{12}\). En
  consecuencia, la función de probabilidad de esta V.A discreta és la
  sigiente: \(\\\) \(\\\) \[P_{X}(x)= \left\{ \begin{array}{lcc}
             \frac{1}{12} & si & x \in [1,12]
             \\
            \\0 &   si  & x= \mbox{cualquier otro número}
            \end{array}
  \right.\]
\end{enumerate}

También nos piden que calculemos la la Esperanza y la Varianza de X.
Para calcular la esoeranza de X, sabemos que
\(E(X)= \sum x\cdot P_{X}(x)\) para toda \(x \in X(\Omega)\). Entonces
tenemos que: \(\\\)
\[E(X)=1\cdot P(1)+2\cdot P(2)+3\cdot P(3)+4\cdot P(4)+5\cdot P(5)+6\cdot P(6)+7\cdot P(7)+8\cdot P(8)+9\cdot P(9)+10\cdot P(10)+11\cdot P(11)+12\cdot P(12)=6.5\]
\[\\\] En cuanto a la varianza, sabemos que
\(Var(X)= E((x-E(X))^2)= E((x-6.5)^2)\). En consecuencia, la varianza de
X nos quedará de la siguiente manera:
\[Var(X) = \sum (x-6.5)^2 \cdot P_{X}(x) = \frac{143}{12}= 11.91666\]
\[\\\]

\begin{enumerate}
\def\labelenumi{\arabic{enumi}.}
\setcounter{enumi}{1}
\tightlist
\item
  La función de distribución de la V.A \(X\) es la siguiente: \(\\\)
  \(\\\) \[F_{X}(x)= \left\{ \begin{array}{lcc}
               0  & si & x<1
          \\
           \\ \frac{1}{12} & si & x=1
           \\
            \\ \frac{2}{12} & si & x=2
            \\
            \\ \frac{3}{12} &   si  & x=3
            \\
            \\ \frac{4}{12} &   si  & x=4
            \\
            \\ \frac{5}{12} &   si  & x=5
            \\
            \\ \frac{6}{12} &   si  & x=6
            \\
            \\ \frac{7}{12} &   si  & x=7
            \\
            \\ \frac{8}{12} &   si  & x=8
            \\
            \\ \frac{9}{12} &   si  & x=9
            \\
            \\ \frac{10}{12} &   si  & x=10
            \\
            \\ \frac{11}{12} &   si  & x=11
            \\
            \\ 1 &   si  & x \ge 12
            \end{array}
  \right.\]
\end{enumerate}

En cuanto al cuantil de \(X\) \(0.4\), basta que observemos la función
de distribución que acabamos de hacer en el apartado anterior. El primer
valor de x que toma la V.A que es mayor o igual que 0.4 es \(x_{0.4}=5\)
ya que \(P(X \le 5)= \frac{5}{12}= 0.4166667\)

\begin{enumerate}
\def\labelenumi{\arabic{enumi}.}
\setcounter{enumi}{2}
\tightlist
\item
  La V.A Y tiene una distribución notable geométrica. Por lo tanto, su
  función de distribución es la siguiente: \[\\\]
  \[F_{Y}(y)= P(Y \le y)=\left\{ \begin{array}{lcc}
               0  & si & y<1
          \\
           \\ 1-(1-\frac{1}{12})^{k+1} & si & \left\{ \begin{array}{lcc}
               k \le x < k+1
           \\ para \mbox{ } k=0,1,2,3...
            \end{array}
  \right.
            \end{array}
  \right.\]
\end{enumerate}

Donde k representa el numero de fallos hasta conseguir el exito \(= 5\)
y \(p= \frac{1}{12}\) es la probabilidad de conseguir el éxito.

\(\\\) 4. Usando la formula de la esperanza y de la varianza de las
distribuciones notables geométricas, obtememos que:
\[E(Y)= \frac{1-p}{p}=\frac{1-\frac{1}{12}}{\frac{1}{12}}= 11\] \[\\\]
\[Var(Y)= \frac{1-p}{p^2}=\frac{1-\frac{1}{12}}{\frac{1}{144}}= 132\]
\[\\\]

\hypertarget{soluciuxf3n-con-r-ejercicio-4}{%
\subsection{\texorpdfstring{Solución con \texttt{R} Ejercicio
4}{Solución con R Ejercicio 4}}\label{soluciuxf3n-con-r-ejercicio-4}}

\begin{Shaded}
\begin{Highlighting}[]
\CommentTok{\#Apartado 1}

\NormalTok{p}\OtherTok{=}\DecValTok{1}\SpecialCharTok{/}\DecValTok{12}
\NormalTok{p}
\end{Highlighting}
\end{Shaded}

\begin{verbatim}
## [1] 0.08333333
\end{verbatim}

\begin{Shaded}
\begin{Highlighting}[]
\NormalTok{Esperanza3}\OtherTok{=}\NormalTok{(}\DecValTok{1}\SpecialCharTok{/}\DecValTok{12}\NormalTok{)}\SpecialCharTok{+}\NormalTok{(}\DecValTok{2}\SpecialCharTok{/}\DecValTok{12}\NormalTok{)}\SpecialCharTok{+}\NormalTok{(}\DecValTok{3}\SpecialCharTok{/}\DecValTok{12}\NormalTok{)}\SpecialCharTok{+}\NormalTok{(}\DecValTok{4}\SpecialCharTok{/}\DecValTok{12}\NormalTok{)}\SpecialCharTok{+}\NormalTok{(}\DecValTok{5}\SpecialCharTok{/}\DecValTok{12}\NormalTok{)}\SpecialCharTok{+}\NormalTok{(}\DecValTok{6}\SpecialCharTok{/}\DecValTok{12}\NormalTok{)}\SpecialCharTok{+}\NormalTok{(}\DecValTok{7}\SpecialCharTok{/}\DecValTok{12}\NormalTok{)}\SpecialCharTok{+}\NormalTok{(}\DecValTok{8}\SpecialCharTok{/}\DecValTok{12}\NormalTok{)}\SpecialCharTok{+}\NormalTok{(}\DecValTok{9}\SpecialCharTok{/}\DecValTok{12}\NormalTok{)}\SpecialCharTok{+}\NormalTok{(}\DecValTok{10}\SpecialCharTok{/}\DecValTok{12}\NormalTok{)}\SpecialCharTok{+}\NormalTok{(}\DecValTok{11}\SpecialCharTok{/}\DecValTok{12}\NormalTok{)}\SpecialCharTok{+}\NormalTok{(}\DecValTok{12}\SpecialCharTok{/}\DecValTok{12}\NormalTok{)}
\NormalTok{Esperanza3}
\end{Highlighting}
\end{Shaded}

\begin{verbatim}
## [1] 6.5
\end{verbatim}

\begin{Shaded}
\begin{Highlighting}[]
\NormalTok{Varianza3}\OtherTok{=}\NormalTok{(}\DecValTok{1}\SpecialCharTok{/}\DecValTok{12}\NormalTok{)}\SpecialCharTok{*}\NormalTok{(((}\DecValTok{1}\FloatTok{{-}6.5}\NormalTok{)}\SpecialCharTok{\^{}}\DecValTok{2}\NormalTok{)}\SpecialCharTok{+}\NormalTok{((}\DecValTok{2}\FloatTok{{-}6.5}\NormalTok{)}\SpecialCharTok{\^{}}\DecValTok{2}\NormalTok{)}\SpecialCharTok{+}\NormalTok{((}\DecValTok{3}\FloatTok{{-}6.5}\NormalTok{)}\SpecialCharTok{\^{}}\DecValTok{2}\NormalTok{)}\SpecialCharTok{+}\NormalTok{((}\DecValTok{4}\FloatTok{{-}6.5}\NormalTok{)}\SpecialCharTok{\^{}}\DecValTok{2}\NormalTok{)}\SpecialCharTok{+}\NormalTok{((}\DecValTok{5}\FloatTok{{-}6.5}\NormalTok{)}\SpecialCharTok{\^{}}\DecValTok{2}\NormalTok{)}\SpecialCharTok{+}\NormalTok{((}\DecValTok{6}\FloatTok{{-}6.5}\NormalTok{)}\SpecialCharTok{\^{}}\DecValTok{2}\NormalTok{)}\SpecialCharTok{+}
\NormalTok{                    ((}\DecValTok{7}\FloatTok{{-}6.5}\NormalTok{)}\SpecialCharTok{\^{}}\DecValTok{2}\NormalTok{)}\SpecialCharTok{+}\NormalTok{((}\DecValTok{8}\FloatTok{{-}6.5}\NormalTok{)}\SpecialCharTok{\^{}}\DecValTok{2}\NormalTok{)}\SpecialCharTok{+}\NormalTok{((}\DecValTok{9}\FloatTok{{-}6.5}\NormalTok{)}\SpecialCharTok{\^{}}\DecValTok{2}\NormalTok{)}\SpecialCharTok{+}\NormalTok{((}\DecValTok{10}\FloatTok{{-}6.5}\NormalTok{)}\SpecialCharTok{\^{}}\DecValTok{2}\NormalTok{)}\SpecialCharTok{+}\NormalTok{((}\DecValTok{11}\FloatTok{{-}6.5}\NormalTok{)}\SpecialCharTok{\^{}}\DecValTok{2}\NormalTok{)}\SpecialCharTok{+}\NormalTok{((}\DecValTok{12}\FloatTok{{-}6.5}\NormalTok{)}\SpecialCharTok{\^{}}\DecValTok{2}\NormalTok{))}
\NormalTok{Varianza3}
\end{Highlighting}
\end{Shaded}

\begin{verbatim}
## [1] 11.91667
\end{verbatim}

\begin{Shaded}
\begin{Highlighting}[]
\CommentTok{\#Apartado 2}

\NormalTok{Cuantil2}\OtherTok{=}\FunctionTok{qgeom}\NormalTok{(}\FloatTok{0.4}\NormalTok{,}\DecValTok{1}\SpecialCharTok{/}\DecValTok{12}\NormalTok{)}
\NormalTok{Cuantil2}
\end{Highlighting}
\end{Shaded}

\begin{verbatim}
## [1] 5
\end{verbatim}

\begin{Shaded}
\begin{Highlighting}[]
\CommentTok{\#Apartado 3}

\CommentTok{\#Vamos a calcular la función de distribución para unos valores de k=5, k=7 y k=9 intentos}
\NormalTok{dist1}\OtherTok{=} \FunctionTok{pgeom}\NormalTok{(}\DecValTok{5}\NormalTok{,}\DecValTok{1}\SpecialCharTok{/}\DecValTok{12}\NormalTok{)}
\NormalTok{dist1}
\end{Highlighting}
\end{Shaded}

\begin{verbatim}
## [1] 0.4067078
\end{verbatim}

\begin{Shaded}
\begin{Highlighting}[]
\NormalTok{dist2}\OtherTok{=} \FunctionTok{pgeom}\NormalTok{(}\DecValTok{7}\NormalTok{,}\DecValTok{1}\SpecialCharTok{/}\DecValTok{12}\NormalTok{)}
\NormalTok{dist2}
\end{Highlighting}
\end{Shaded}

\begin{verbatim}
## [1] 0.5014698
\end{verbatim}

\begin{Shaded}
\begin{Highlighting}[]
\NormalTok{dist3}\OtherTok{=} \FunctionTok{pgeom}\NormalTok{(}\DecValTok{9}\NormalTok{,}\DecValTok{1}\SpecialCharTok{/}\DecValTok{12}\NormalTok{)}
\NormalTok{dist3}
\end{Highlighting}
\end{Shaded}

\begin{verbatim}
## [1] 0.5810961
\end{verbatim}

\begin{Shaded}
\begin{Highlighting}[]
\CommentTok{\#Apartado 4}
\NormalTok{Esperanza4}\OtherTok{=}\NormalTok{((}\DecValTok{1}\SpecialCharTok{{-}}\NormalTok{(}\DecValTok{1}\SpecialCharTok{/}\DecValTok{12}\NormalTok{)))}\SpecialCharTok{/}\NormalTok{(}\DecValTok{1}\SpecialCharTok{/}\DecValTok{12}\NormalTok{)}
\NormalTok{Esperanza4}
\end{Highlighting}
\end{Shaded}

\begin{verbatim}
## [1] 11
\end{verbatim}

\begin{Shaded}
\begin{Highlighting}[]
\NormalTok{Varianza4}\OtherTok{=}\NormalTok{((}\DecValTok{1}\SpecialCharTok{{-}}\NormalTok{(}\DecValTok{1}\SpecialCharTok{/}\DecValTok{12}\NormalTok{)))}\SpecialCharTok{/}\NormalTok{(}\DecValTok{1}\SpecialCharTok{/}\DecValTok{144}\NormalTok{)}
\NormalTok{Varianza4}
\end{Highlighting}
\end{Shaded}

\begin{verbatim}
## [1] 132
\end{verbatim}

\hypertarget{problema-5}{%
\subsection{Problema 5}\label{problema-5}}

La proporción de niños pelirrojos es 1 cada 100. En una ciudad se
produjeron 500 nacimientos (independientes) nacimientos en 2020, modelad
mediante una distribución binomial la variable \(X\) = número de niños
pelirrojos nacidos entre los 500 niños. Utilizad R para calcular de
forma exacta

\begin{enumerate}
\def\labelenumi{\arabic{enumi}.}
\tightlist
\item
  La probabilidad de que ninguno de los nacidos ese año sea pelirrojo.
\item
  La probabilidad de que nazcan más de 2 niños pelirrojos
\item
  Repetir los cálculos con R utilizando una aproximación Poisson
\end{enumerate}

\hypertarget{soluciuxf3n-ejercicio-5}{%
\subsection{Solución Ejercicio 5}\label{soluciuxf3n-ejercicio-5}}

\begin{enumerate}
\def\labelenumi{\arabic{enumi}.}
\tightlist
\item
  Sabiendo que X representa el número de pelirrojos nacidos entre los
  500 niños manejaremos una \(Binom(n=500,p=\frac{1}{100}=0.01)\)
\end{enumerate}

\[P_{X}(x)=\binom{n}{x}\cdot(p^x)+(1-p)^{(n-x)}\]
\[P_{X}(0)=\binom{500}{0}\cdot(p^0)\cdot(1-0.01)^{(500-0)}=0.006570483\]
\[\\\] 2. \[F_{X}(x)=P(X\leq X)\] \[\\\] \[P(X>2)=1-[P(X\leq 2)]\]

\[P(X\leq2)=P_{X}(0)+P_{X}(1)+P_{X}(2)\] \[\\\] \[P(X\leq2)=0.1233858\]
\[\\\] \[P(X>2)=1-[0.1233858]\] \[\\\] \[P(X>2)=0.8766142\] \[\\\] 3.
Sabemos que el límite de una serie de variables
\(B(n,p_n=\frac{\lambda}{n})\). Por lo tanto, nuestra distribución
binomial puede llevarse a cabo mediante un proceso
\(Po(\lambda=0.01\cdot500=5)\). \[\\\]
\[ P(X=0)=\frac{\lambda^x}{x!}\cdot e^{-\lambda}=\frac{5^0}{0!}\cdot e^{-5}=\frac{1}{e^{5}}=0.0067379 \]
\[\\\] \[P(X>2)=1-P(X\leq 2)=1-0.0124652=0.875348\] \[\\\]
\[P(X\leq 2)=P(0)+P(1)+P(2)=0.124652\]

\hypertarget{soluciuxf3n-con-r-ejercicio-5}{%
\subsection{\texorpdfstring{Solución con \texttt{R} Ejercicio
5}{Solución con R Ejercicio 5}}\label{soluciuxf3n-con-r-ejercicio-5}}

\begin{Shaded}
\begin{Highlighting}[]
\DocumentationTok{\#\#Apartado 1}

\DocumentationTok{\#\#Con la función dbinom podemos calcular exactamente la función de densidad de }
\DocumentationTok{\#\#v.a discreta binomial para un determinado valor de X. Por parámetros damos el }
\DocumentationTok{\#\#valor de n y la probabilidad correspondiente.}
\FunctionTok{dbinom}\NormalTok{(}\DecValTok{0}\NormalTok{,}\AttributeTok{size=}\DecValTok{500}\NormalTok{,}\AttributeTok{prob=}\FloatTok{0.01}\NormalTok{)}
\end{Highlighting}
\end{Shaded}

\begin{verbatim}
## [1] 0.006570483
\end{verbatim}

\begin{Shaded}
\begin{Highlighting}[]
\DocumentationTok{\#\#Otra posibilidad es ir haciendo los cálculos correspondientes a la fórmula:}
\NormalTok{y}\OtherTok{=}\FunctionTok{choose}\NormalTok{(}\DecValTok{500}\NormalTok{,}\DecValTok{0}\NormalTok{)}\SpecialCharTok{*}\NormalTok{(}\FloatTok{0.01}\SpecialCharTok{\^{}}\DecValTok{0}\NormalTok{)}\SpecialCharTok{*}\NormalTok{(}\DecValTok{1}\FloatTok{{-}0.01}\NormalTok{)}\SpecialCharTok{\^{}}\NormalTok{(}\DecValTok{500{-}0}\NormalTok{)}
\NormalTok{y}
\end{Highlighting}
\end{Shaded}

\begin{verbatim}
## [1] 0.006570483
\end{verbatim}

\begin{Shaded}
\begin{Highlighting}[]
\DocumentationTok{\#\#Apartado 2}

\DocumentationTok{\#\#Con la función pbinom podemos calcular la función de probabilidad de una }
\DocumentationTok{\#\#v.a discreta binomial para un intervalo superior a determinado X. }
\DocumentationTok{\#\#Por parámetros damos el valor de X, la n, la probabilidad y el lower.tail}
\DocumentationTok{\#\#(nos indica si tenemos \textless{}= o  \textgreater{}=).}
\FunctionTok{pbinom}\NormalTok{(}\DecValTok{2}\NormalTok{,}\AttributeTok{size=}\DecValTok{500}\NormalTok{,}\AttributeTok{prob=}\FloatTok{0.01}\NormalTok{,}\AttributeTok{lower.tail=}\ConstantTok{FALSE}\NormalTok{)}
\end{Highlighting}
\end{Shaded}

\begin{verbatim}
## [1] 0.8766142
\end{verbatim}

\begin{Shaded}
\begin{Highlighting}[]
\DocumentationTok{\#\#Apartado 3}

\DocumentationTok{\#\#Con la función dpois podemos calcular la función de densidad de una v.a}
\DocumentationTok{\#\# discreta de Poisson.Parámetros: valor de la x y lambda.}
\FunctionTok{dpois}\NormalTok{(}\DecValTok{0}\NormalTok{,}\DecValTok{5}\NormalTok{)}
\end{Highlighting}
\end{Shaded}

\begin{verbatim}
## [1] 0.006737947
\end{verbatim}

\begin{Shaded}
\begin{Highlighting}[]
\DocumentationTok{\#\#Sin embargo, con la función ppois ppodemos calcular la función de }
\DocumentationTok{\#\#distribución.Parámetros: x,lambda y orientación del signo \textless{}/\textgreater{}}
\FunctionTok{ppois}\NormalTok{(}\DecValTok{2}\NormalTok{,}\DecValTok{5}\NormalTok{,}\AttributeTok{lower.tail=}\ConstantTok{FALSE}\NormalTok{)}
\end{Highlighting}
\end{Shaded}

\begin{verbatim}
## [1] 0.875348
\end{verbatim}

\hypertarget{problema-6}{%
\subsection{Problema 6}\label{problema-6}}

Las consultas a una base datos llegan a un ritmo de medio \(\lambda=5\)
peticiones por segundo. Sabemos que el nombre de peticiones que llegan
en un segundo es una variable aleatoria que aproximadamente tienen una
distribución de Poisson.

\begin{enumerate}
\def\labelenumi{\arabic{enumi}.}
\tightlist
\item
  Calcular la probabilidad que lleguen más de 10 peticiones en 3
  segundos utilizad R.
\item
  Calcular que entre una consulta y la siguiente pasen más de \(0.5\)
  segundos.
\item
  Calcular el cuantil 0.5 de \(X_{t=10}\) numero de peticiones en 10
  segundos utilizad R.
\end{enumerate}

\hypertarget{soluciuxf3n-ejercicio-6}{%
\subsection{Solución Ejercicio 6}\label{soluciuxf3n-ejercicio-6}}

En este problema, vamos a trabajar con una distribución
\(Po((\lambda= 5)\cdot t)\), es decir, una distribución Poisson por una
unidad de tiempo (en este caso, el numero de peticiones que llegan por
unidad de tiempo).

\begin{enumerate}
\def\labelenumi{\arabic{enumi}.}
\tightlist
\item
  Sabiendo quen la función de probabilidad de este tipo de V.A es la
  siguiente: \[\\\] \[F_{Y}(y)= P(Y \le y)=\left\{ \begin{array}{lcc}
               \frac{(\lambda \cdot t)^x}{x!} \cdot e^{-\lambda \cdot t}  & si & x=0,1,2,3...
          \\
           \\ 0 && en \mbox{ } otro \mbox{ } caso
            \end{array}
  \right.\]
\end{enumerate}

deberemos llevar a cabo el siguente calculo para encontrar la
probabilidad de que lleguen más de 10 peticiones en 3 segundos: \[\\\]
\[P_{X}(10)=\frac{(5\cdot 3)^{10}}{10!} \cdot e^{-5 \cdot 3}=0.04861075083 \]
\[\\\] 2. Sea T una V.A que contabiliza el tiempo transcurrido entre dos
peticiones consecutivas y considerando que T solo puede tomar valores
positivos \((t>0)\), entonces podemos decir que :
\(P(T>t=0.5)= P(\mbox{cero eventos en el intervalo } (0,0.5])\). Por
tanto tenemos que : \[\\\]
\[P_{X}(x=0)= \frac{(5\cdot 0.5)^{0}}{0!}\cdot e^{-5\cdot 0.5}= 0.08208499862\]
\[\\\] 3. Si queremos calcular el cuantil de 0.5 cuando se da el numero
de peticiones por cada 10 segundos, queremos calcular lo siguente:
\[\\\] \[P(X \le x_{0.5}) \ge 0.5\] \[\\\] Para ello, usaremos el
siguente comando en la consola de \texttt{R}:
\(qpois(0.5, \lambda \cdot t=5 \cdot 10=50)\). (Ver resultados en el
Chunk "Ejercicio 6) \[\\\]

\hypertarget{soluciuxf3n-con-r-ejercicio-6}{%
\subsection{\texorpdfstring{Solución con \texttt{R} Ejercicio
6}{Solución con R Ejercicio 6}}\label{soluciuxf3n-con-r-ejercicio-6}}

\begin{Shaded}
\begin{Highlighting}[]
\CommentTok{\#Apartado 1}

\NormalTok{P10}\OtherTok{=}\FunctionTok{dpois}\NormalTok{(}\DecValTok{10}\NormalTok{,}\DecValTok{5}\SpecialCharTok{*}\DecValTok{3}\NormalTok{)}
\NormalTok{P10}
\end{Highlighting}
\end{Shaded}

\begin{verbatim}
## [1] 0.04861075
\end{verbatim}

\begin{Shaded}
\begin{Highlighting}[]
\CommentTok{\#Apartado 2}

\NormalTok{P0}\OtherTok{=}\FunctionTok{dpois}\NormalTok{(}\DecValTok{0}\NormalTok{,}\DecValTok{5}\SpecialCharTok{*}\FloatTok{0.5}\NormalTok{)}
\NormalTok{P0}
\end{Highlighting}
\end{Shaded}

\begin{verbatim}
## [1] 0.082085
\end{verbatim}

\begin{Shaded}
\begin{Highlighting}[]
\CommentTok{\#Apartado 3}

\NormalTok{Cuantil3}\OtherTok{=}\FunctionTok{qpois}\NormalTok{(}\FloatTok{0.5}\NormalTok{, }\DecValTok{5}\SpecialCharTok{*}\DecValTok{10}\NormalTok{)}
\NormalTok{Cuantil3}
\end{Highlighting}
\end{Shaded}

\begin{verbatim}
## [1] 50
\end{verbatim}

\hypertarget{problema-7}{%
\subsection{Problema 7}\label{problema-7}}

Tenemos que elegir entre dos programas (Prog1 y Prog2), el objetivo es
elegir el programa más rápido en tiempo de respuesta en nuestro cluster
de ordenadores. El tiempo de ejecución del Prog1 se ha modelado según
una \(N(\mu_1=100, \sigma_1=300)\) (la probabilidad de un tiempo de
ejecución negativo es despreciable) y en Prog2 según una
\(N(\mu_2=90, \sigma_2=300)\). Utilizad R para el cálculo final de las
probabilidades de la normal. (Utilizad R para el cálculo final)

\begin{enumerate}
\def\labelenumi{\arabic{enumi}.}
\tightlist
\item
  ¿Qué Programa elegimos si queremos que el el 90\% de los casos el
  tiempo de respuesta sea menor ?
\item
  Calcular la probabilidad de que el tiempo de ejecución sea mayor que
  130 para cada algoritmo.
\end{enumerate}

\hypertarget{soluciuxf3n-ejercicio-7}{%
\subsection{Solución Ejercicio 7}\label{soluciuxf3n-ejercicio-7}}

\begin{enumerate}
\def\labelenumi{\arabic{enumi}.}
\tightlist
\item
  Para saber cual tiene una respuesta más rápida el \(90%
  \) de las veces, llevaremos a cabo el calculo del cuantil 0.9 para
  ambas variables:
  \[x_{0.9}=qnorm(0.9, mean=100, sd=300)=484.4655 \Rightarrow Prog1\]
  \[\\\]
  \[x_{0.9}=qnorm(0.9, mean=90, sd=300)=474.4655 \Rightarrow Prog2\]
  \[\\\] Viendo los siguentes resultados, el tiempo de respuesta de el
  programa 2 es menor. Por ese motivo, elegiremos dicho programa. \[\\\]
\item
  Para realizar este calculo, deberemos llevar a cabo la tipicación de
  las variables aleatorias \(X_{1} \mbox{ y } X_{2}\). Para ello,
  realizaremos la siguente transformacion lineal a ambas V.A: \[\\\]
  \[Z_{1}=\frac{X_{1}-\mu_{1}}{\sigma_{1}}\] \[\\\]
  \[Z_{2}=\frac{X_{2}-\mu_{2}}{\sigma_{2}}\] \[\\\] De esta manera,
  pasaremos a trabajar con dos distribuciones normales estandar
  \(N_{1}(0,1)\) y \(N_{2}(0,1)\).Teniendo esto en cuenta, las
  probabilidades que nos piden que calculemos son als siguentes:
  \[P(Z_{1}>\frac{130-100}{300})=1-P(Z_{1}\le \frac{1}{10})=1-F_{Z_{1}}(\frac{1}{10})=1-0.5398278=0.4601722\]
  \[\\\]
  \[P(Z_{2}>\frac{130-90}{300})=1-F_{Z_{2}}(\frac{2}{15})=1-P(Z_{2}\le \frac{2}{15})=1-0.5530351=0.4469649\]
  \[\\\]
\end{enumerate}

\hypertarget{soluciuxf3n-con-r-ejercicio-7}{%
\subsection{\texorpdfstring{Solución con \texttt{R} Ejercicio
7}{Solución con R Ejercicio 7}}\label{soluciuxf3n-con-r-ejercicio-7}}

\begin{Shaded}
\begin{Highlighting}[]
\CommentTok{\#Apartado 1}

\NormalTok{cuantilProg1}\OtherTok{=}\FunctionTok{qnorm}\NormalTok{(}\FloatTok{0.9}\NormalTok{, }\DecValTok{100}\NormalTok{, }\DecValTok{300}\NormalTok{)}
\NormalTok{cuantilProg1}
\end{Highlighting}
\end{Shaded}

\begin{verbatim}
## [1] 484.4655
\end{verbatim}

\begin{Shaded}
\begin{Highlighting}[]
\NormalTok{cuantilProg2}\OtherTok{=}\FunctionTok{qnorm}\NormalTok{(}\FloatTok{0.9}\NormalTok{, }\DecValTok{90}\NormalTok{, }\DecValTok{300}\NormalTok{)}
\NormalTok{cuantilProg2}
\end{Highlighting}
\end{Shaded}

\begin{verbatim}
## [1] 474.4655
\end{verbatim}

\begin{Shaded}
\begin{Highlighting}[]
\CommentTok{\#Apartado 2}

\CommentTok{\#No haria falta poner lo parametros de mu i lambda, ya que vienen por defecto como mean=0 y sd=1}
\CommentTok{\#Pero los pongo para que quede más claro}
\NormalTok{PZ1}\OtherTok{=}\FunctionTok{pnorm}\NormalTok{(}\DecValTok{1}\SpecialCharTok{/}\DecValTok{10}\NormalTok{,}\DecValTok{0}\NormalTok{,}\DecValTok{1}\NormalTok{)}
\NormalTok{PZ1}
\end{Highlighting}
\end{Shaded}

\begin{verbatim}
## [1] 0.5398278
\end{verbatim}

\begin{Shaded}
\begin{Highlighting}[]
\NormalTok{PZ2}\OtherTok{=}\FunctionTok{pnorm}\NormalTok{(}\DecValTok{2}\SpecialCharTok{/}\DecValTok{15}\NormalTok{,}\DecValTok{0}\NormalTok{,}\DecValTok{1}\NormalTok{)}
\NormalTok{PZ2}
\end{Highlighting}
\end{Shaded}

\begin{verbatim}
## [1] 0.5530351
\end{verbatim}

\hypertarget{problema-8}{%
\subsection{Problema 8}\label{problema-8}}

En la NBA el
\href{https://es.wikipedia.org/wiki/Jos\%C3\%A9_Manuel_Calder\%C3\%B3n}{José
Calderón} fue en la temporada
\href{https://es.wikipedia.org/wiki/L\%C3\%ADderes_en_porcentaje_de_tiros_libres_de_la_NBA}{2008-09
el jugador de baloncesto} con mejor porcentaje tiros libres anotados un
98.05\%.

Justificar los cálculos con notación matemática y haced el cálculo final
con R.

\begin{enumerate}
\def\labelenumi{\arabic{enumi}.}
\tightlist
\item
  ¿Cual es el valor esperando y la varianza del número tiros hasta
  acertar 10 tiros libres?
\item
  ¿Cuál es la probabilidad de que acierte al menos 40 tiros libres de
  forma consecutiva.
\item
  ¿Cuál es la probabilidad de que haga una serie de 100 tiros hasta
  obtener el tercer fallo?
\end{enumerate}

\hypertarget{soluciuxf3n-ejercicio-8}{%
\subsection{Solución Ejercicio 8}\label{soluciuxf3n-ejercicio-8}}

\begin{enumerate}
\def\labelenumi{\arabic{enumi}.}
\tightlist
\item
  Consideramos un esperimento bernoulli con probabilidad de éxito
  \(p_1=0.9805\). Dado este experimento, podemos considerar una V.A
  \(X_1\) de distribución Binomial Negativa, que nos cuenta el numero de
  tiros libres fallados hasta que acertamos 10. Por lo tanto,
  trabajaremos con una \(BN(n=10 \mbox{ aciertos }, p=0.9805)\).
  Entonces tenemos que: \[\\\]
  \[E(X_1)=\frac{n\cdot(1-p_1)}{p_1}=\frac{10\cdot(1-0.9805)}{0.9805}=0.1988781\]
  \[\\\]
  \[Var(X_1)=\frac{n\cdot(1-p_{1})}{p_{1}^{2}}=\frac{10\cdot(1-0.9805)}{0.9805^2}=0.2028334\]
  \[\\\]
\item
  Consideremos ahora un experimento Bernoulli con probabilidad de éxito
  \(p_2=1-0.9805=0.0195\) (probabilidad de fallar un tiro libre).
  Consideramos ahora una V.A \(X_2\) que nos cuenta el numero de fallos
  hasta conseguir el primer éxito (éxito=fallar el tiro libre, fracaso=
  acertar el tiro libre). Por lo tanto, trabajaremos con una V.A de tipo
  Geométrica \(\Rightarrow Ge(p_2=0.0195)\). Nos piden: \[\\\]
  \[P(X_{2}>40)= 1-P(X_{2}\le 40)= 1-F_{X_{2}}(40)=1-(1-(1-p_{2})^{k})= 1-(1-(1-0.0195)^{40})=0.4548874134\]
  \[\\\]
\item
  Para este caso, consideramos una V.A \(X\) que nos cuenta el numero de
  éxitos de una cadena de n intentos del experimento aleatorio ``fallar
  un tiro libre'', con probabilidad de éxito \(p_{3}=0.0195\). Entonces,
  nuestra V.A sigue una distribución binomial
  \(B(n=99 \mbox{ lanzamientos }, p_{3}=0.0195)\), ya que sabemos que el
  ultimo lanzamiento siempre será un éxito (por eso 99 y no 100
  lanzamientos). Finalmente, nos piden:
  \[P(X_{3}=3)=\binom{n}{x_{3}} \cdot p_{3}^{x_{3}} \cdot (1-p_{3})^{n-x_3} = \binom{99}{3} \cdot 0.0195^{3} \cdot (1-0.0195)^{99-3}= 0.1756123\]
  \[\\\]
\end{enumerate}

\hypertarget{soluciuxf3n-con-r-ejercicio-8}{%
\subsection{\texorpdfstring{Solución con \texttt{R} Ejercicio
8}{Solución con R Ejercicio 8}}\label{soluciuxf3n-con-r-ejercicio-8}}

\begin{Shaded}
\begin{Highlighting}[]
\CommentTok{\#Apartado 1}

\NormalTok{Esperanza5}\OtherTok{=}\NormalTok{(}\DecValTok{10}\SpecialCharTok{*}\NormalTok{(}\DecValTok{1}\FloatTok{{-}0.9805}\NormalTok{))}\SpecialCharTok{/}\FloatTok{0.9805}
\NormalTok{Esperanza5}
\end{Highlighting}
\end{Shaded}

\begin{verbatim}
## [1] 0.1988781
\end{verbatim}

\begin{Shaded}
\begin{Highlighting}[]
\NormalTok{Varianza5}\OtherTok{=}\NormalTok{(}\DecValTok{10}\SpecialCharTok{*}\NormalTok{(}\DecValTok{1}\FloatTok{{-}0.9805}\NormalTok{))}\SpecialCharTok{/}\NormalTok{(}\FloatTok{0.9805}\SpecialCharTok{\^{}}\DecValTok{2}\NormalTok{)}
\NormalTok{Varianza5}
\end{Highlighting}
\end{Shaded}

\begin{verbatim}
## [1] 0.2028334
\end{verbatim}

\begin{Shaded}
\begin{Highlighting}[]
\CommentTok{\#Apartado 2}

\CommentTok{\#Ponemos el valor de la variable = 39, porque R implementa la función geometrica}
\CommentTok{\#que empieza a contar des del 0.}
\NormalTok{D40}\OtherTok{=}\DecValTok{1}\SpecialCharTok{{-}}\FunctionTok{pgeom}\NormalTok{(}\DecValTok{39}\NormalTok{,}\FloatTok{0.0195}\NormalTok{)}
\NormalTok{D40}
\end{Highlighting}
\end{Shaded}

\begin{verbatim}
## [1] 0.4548874
\end{verbatim}

\begin{Shaded}
\begin{Highlighting}[]
\CommentTok{\#Apartado 3}

\NormalTok{P3}\OtherTok{=}\FunctionTok{dbinom}\NormalTok{(}\DecValTok{3}\NormalTok{, }\DecValTok{99}\NormalTok{, }\FloatTok{0.0195}\NormalTok{)}
\NormalTok{P3}
\end{Highlighting}
\end{Shaded}

\begin{verbatim}
## [1] 0.1756123
\end{verbatim}

\end{document}
